Мета роботи: дослідження принципів побудови та застосування сучасних асиметричних 
криптосистем, отримання практичних навичок використання механізмів цифрового 
підпису та узгодження ключів, аналіз особливостей безпечного застосування 
асиметричних криптографічних алгоритмів. 

\section{Дослідження принципів асиметричної криптографії}
\subsection{У чому полягає основний принцип асиметричної криптографії та які функції виконують відкритий і закритий ключі?}

Основний принцип асиметричної криптографії полягає у використанні пари математично 
пов'язаних ключів: відкритого та закритого. Відкритий ключ можна вільно передавати іншим користувачам, 
тоді як закритий ключ повинен зберігатися в таємниці його власником.

Залежно від призначення криптографічної операції ключі виконують різні функції. 
Під час шифрування повідомлення використовується відкритий ключ одержувача, 
а для його розшифрування — відповідний закритий ключ. У випадку цифрового підпису, навпаки, 
закритий ключ використовується для створення підпису, а відкритий — для його перевірки. 
Таким чином, асиметрична криптографія дозволяє забезпечити конфіденційність інформації, 
автентичність її джерела та цілісність повідомлень \cite{RFC8017}.


\subsection{Які принципові відмінності між асиметричним шифруванням, цифровим підписом та узгодженням спільного секрету (key agreement)?}

Ці три механізми мають різне призначення.

Асиметричне шифрування використовується для забезпечення конфіденційності інформації. 
Відправник шифрує повідомлення відкритим ключем одержувача, а розшифрувати його може 
лише власник відповідного закритого ключа.

Цифровий підпис призначений для підтвердження автентичності відправника та цілісності повідомлення. 
Він створюється за допомогою закритого ключа відправника, а перевіряється за допомогою його відкритого ключа. 
Зміна повідомлення призводить до того, що перевірка підпису завершується невдало.

Узгодження спільного секрету дозволяє двом сторонам, які обмінюються відкритими даними, 
незалежно отримати однаковий секретний ключ. Надалі цей ключ можна використовувати для симетричного шифрування інформації.

Отже, основна відмінність полягає в результаті їх використання: шифрування захищає зміст повідомлення, 
цифровий підпис підтверджує його походження та цілісність, а узгодження ключів дозволяє отримати спільний 
секрет для подальшого захищеного обміну даними \cite{RFC8017,RFC7748}.


\subsection{Які задачі вирішують RSA, ECDSA/EdDSA та ECDH/X25519? Чому ці алгоритми не можна вважати взаємозамінними?}

RSA, ECDSA, EdDSA та ECDH/X25519 є криптографічними алгоритмами з відкритим ключем, проте вони відрізняються за призначенням.

RSA — алгоритм, який може використовуватися як для асиметричного шифрування, так і для створення та перевірки цифрових підписів. 
Його стійкість ґрунтується на складності факторизації великих цілих чисел.

ECDSA та EdDSA призначені для створення і перевірки цифрових підписів на основі криптографії еліптичних кривих. 
ECDSA використовує випадкове значення під час створення кожного підпису, тоді як EdDSA
використовує детермінований механізм формування підпису.

ECDH та X25519 використовуються для узгодження спільного секрету між двома сторонами. Вони дозволяють отримати однаковий секрет, 
не передаючи його безпосередньо мережею.

Ці алгоритми не є взаємозамінними, оскільки розв'язують різні криптографічні задачі. Наприклад, ECDSA та EdDSA забезпечують цифровий підпис, 
але не призначені для узгодження спільного секрету. Натомість ECDH та X25519 дозволяють узгодити секрет, але самі по собі не забезпечують цифрового підпису. 
RSA має ширший набір застосувань, однак це не означає, що він може безпосередньо замінити будь-який інший алгоритм без зміни протоколу та способу його використання \cite{RFC8017,RFC8032,RFC7748}.


\subsection{Чому довжини ключів RSA та криптосистем на еліптичних кривих не можна безпосередньо порівнювати для оцінювання рівня криптографічної стійкості?}

Безпосередньо порівнювати довжини ключів RSA та криптосистем на еліптичних кривих некоректно, оскільки їхня стійкість ґрунтується на різних математичних задачах. 
Для RSA основною є складність факторизації великого числа, а для криптосистем на еліптичних кривих — складність задачі дискретного логарифмування в групі точок еліптичної кривої.

Тому однакова кількість бітів у ключах не означає однаковий рівень захисту. Наприклад, ключ RSA довжиною 3072 біти та ключ на еліптичній кривій з 
порядком базової точки приблизно 256 біт забезпечують орієнтовно однаковий рівень класичної криптографічної стійкості — 128 біт.

Отже, для порівняння таких систем потрібно враховувати не лише довжину ключа, а й складність математичної задачі, на якій ґрунтується їхня безпека. 
Саме тому для оцінювання стійкості використовують поняття рівня криптографічної безпеки (security strength), а не просто кількість бітів у ключі \cite{NIST80057}.


\subsection{Для чого при використанні RSA застосовуються спеціальні схеми доповнення (padding)? Які ризики створює неправильний вибір або використання padding?}

Схема доповнення при використанні RSA застосовується для усунення математичних дефектів алгоритму та
забезпечення необхідного рівня криптографічної стійкості. Застосування padding перетворює детерміноване
шифрування на ймовірнісне за рахунок додавання випадкових байтів, що забезпечує формування нового
шифротексту при кожному повторному шифруванні однакових даних~\cite{RFC8017}. Крім того, схеми
доповнення руйнують гомоморфну мультиплікативну структуру шифру та захищають від атак добування
кореня при малих значеннях публічної експоненти. Неправильний вибір або використання застарілих схем
доповнення, таких як PKCS \#1 v1.5, створює ризики реалізації атак на основі витоку інформації про помилки
доповнення, відомих як атака Блейхенбахера або Padding Oracle. Для усунення цих вразливостей стандарт
NIST SP 800-56B вимагає використання сучасних схем RSA-OAEP для шифрування та RSA-PSS для цифрових
підписів \cite{RFC8017,NIST80056B}.

\subsection{Яку роль виконує nonce при формуванні цифрового підпису ECDSA? Які наслідки може мати повторне використання одного nonce для різних повідомлень?}

Під час створення цифрового підпису ECDSA значення nonce виконує роль одноразового випадкового числа,
яке використовується для обчислення тимчасової точки на еліптичній кривій та формування першої
компоненти підпису. Це значення має бути строго унікальним та непередбачуваним для кожного нового
підпису~\cite{RFC6979}. Повторне використання одного й того самого значення nonce для підписання двох
різних повідомлень на одному закритому ключі призводить до критичних наслідків у вигляді повного
математичного розкриття приватного ключа підписувача. Атакуючий може обчислити різницю підписів,
виявити значення secret nonce та на його основі легко відновити закритий ключ. Для запобігання цьому варто
застосовувати детерміновану генерацію nonce або використовувати алгоритм EdDSA, який за своєю
конструкцією усуває ризик повторного використання тимчасового
параметра~\cite{RFC6979,RFC8032,cryptography}.

\subsection{Чому асиметричне шифрування зазвичай не використовується безпосередньо для захисту великих обсягів даних? Що таке гібридна криптосистема?}

Асиметричне шифрування не застосовується безпосередньо для захисту великих обсягів даних через низьку
обчислювальну швидкість алгоритмів та жорсткі обмеження на розмір вхідного блоку. Операції піднесення
до степеня за модулем великих чисел потребують значно більше ресурсів процесора порівняно із
симетричними блоковими шифрами~\cite{NIST80057}. Гібридна криптосистема є комбінованим
криптографічним механізмом, який поєднує переваги асиметричної та симетричної криптографії. У такій
системі сам обсяг даних шифрується швидким симетричним шифром, наприклад AES-256-GCM, за допомогою
одноразового випадкового сеансового ключа. Сам сеансовий ключ має малий розмір і шифрується відкритим
асиметричним ключем отримувача за допомогою RSA-OAEP або механізмів інкапсуляції ключів
KEM~\cite{NIST80057,RFC8017}.


\subsection{Які вимоги висуваються до зберігання та використання закритого ключа асиметричної криптосистеми?}

Вимоги до зберігання та використання закритого ключа є забезпечення його конфіденційності та цілісності на
всіх етапах життєвого циклу~\cite{NIST80057,FIPS1403}. Закритий ключ повинен зберігатися в ізольованому
середовищі із застосуванням суворого розмежування прав доступу. При збереженні на носіях ключ
обов'язково захищається симетричним шифруванням із використанням функцій виведення ключів на основі
паролів, таких як PBKDF2 або Argon2id. Для систем із високим рівнем безпеки вимагається використання
апаратних модулів безпеки HSM, TPM або смарт-карт, які унеможливлюють експорт ключа у відкритому
вигляді за межі пристрою~\cite{FIPS1403}. Під час використання в оперативній пам'яті ключ має захищатися
від витоків через побічні канали, а після завершення операції підлягає негайному та безповоротному
знищенню процедурою zeroization~\cite{NIST80057,FIPS1403}.
